% ── sections/06_limitations.tex ───────────────────────────────────────────────

\section{Limitations}

Several limitations of this study should be noted.

\textbf{Single-seed best result.} The M3 six-network best result
(AUPR~=~0.8067) is reported from a single training run and lacks multi-seed
variance estimation. While the baseline configuration is validated across five
seeds (0.8019~$\pm$~0.0050), the optimal configuration identified through
hyperparameter search has not been similarly characterised.

\textbf{Interpretability model mismatch.} The Integrated Gradients attribution
and gene set enrichment analyses (Section~4.4) were performed on a six-network
EMGNNImproved model, but not on the best heterophily-aware model
(AUPR~=~0.824). Re-running these analyses on the strongest model would provide
a stronger biological validation, but is outside the final experimental scope.

\textbf{Transductive target-network input.} CPDB topology and node features are
included as input while CPDB labels define the test set, following the original
transductive EMGNN protocol. Test labels are withheld, but the high learned CPDB
weight (0.210) may partly reflect access to target-network structure. A stricter
leave-target-network-out evaluation is needed before interpreting the weights as
intrinsic database quality.

\textbf{Legacy split provenance.} The experiments reported here predate the
gene-level JSON split manifest introduced during the reproducibility audit.
Although train and test membership originated from the benchmark masks, the
historical validation subset was not saved independently. The current submission
therefore reports these comparisons as explicitly labelled legacy evidence rather
than as fixed-split revalidation results.

\textbf{Zero-baseline IG attribution.} Integrated Gradients uses a zero-vector
baseline for feature attribution. Because the reported experiments use aligned
raw benchmark features without additional Z-score standardisation, zero does not
have one uniform biological interpretation across all modalities. Baseline
sensitivity analysis was not performed.

\textbf{Out-of-scope extensions.} Four implemented techniques---pathway hypergraph
integration, HIPGNN anomaly detection, PINNACLE pretrained embeddings, and
GNNExplainer---await external dataset preprocessing or further debugging. They are
retained as optional future work and are not used to support the conclusions of
this paper.

\textbf{Comparison scope.} Direct numerical comparison with recent methods
(SGCD, DISHyper, DGHNN, deepCDG) is limited by differences in evaluation
protocols, dataset versions, and train/test splits. Standardised cross-method
benchmarks on identical data splits would enable more rigorous comparison.
